\documentclass[12pt,aspectratio=169]{beamer}
\usetheme{Copenhagen}
\usecolortheme{whale}
\setbeamertemplate{navigation symbols}{}
\setbeamertemplate{footline}[frame number]
\usepackage{amsmath,amssymb}
\usepackage{tikz}
\usepackage{booktabs}
\usepackage{listings}
\usepackage{xeCJK}
\setCJKmainfont{Noto Serif CJK SC}
\setCJKsansfont{Noto Sans CJK SC}
\setCJKmonofont{Noto Sans CJK SC}

\lstset{
  basicstyle=\ttfamily\footnotesize,
  keywordstyle=\color{blue}\bfseries,
  commentstyle=\color{gray},
  stringstyle=\color{red},
  showstringspaces=false,
  numbers=left,
  numberstyle=\tiny\color{gray},
  frame=single,
  tabsize=4,
  language=C
}

\title{稀疏矩阵与 CRS 存储格式}
\subtitle{存储格式 \textperiodcentered\ CRS (Compressed Row Storage) \textperiodcentered\ C 语言实现}
\author{计算方法教学}
\date{\today}

\begin{document}

% ==================== frame 1: 标题 ====================
\begin{frame}
\titlepage
\end{frame}

% ==================== frame 2: 什么是稀疏矩阵 ====================
\begin{frame}{什么是稀疏矩阵？}
\framesubtitle{定义 \textperiodcentered\ 来源 \textperiodcentered\ 存储动机}

\textbf{定义：} 非零元素个数远小于总元素个数的矩阵。\\
设 $N$ 为矩阵总元素数，$N_{\!nz}$ 为非零元个数，若 $N_{\!nz} \ll N$，则称为稀疏矩阵。

\vspace{8pt}
\textbf{常见来源：}
\begin{itemize}
  \item 有限差分 / 有限元离散 PDE — 每个节点只与相邻节点耦合
  \item 图论问题 — 邻接矩阵，每行平均度 $\ll n$
  \item 电路仿真 — 节点导纳矩阵，每个器件只连接少量节点
  \item PageRank — 网页链接矩阵，平均出度 $O(1)$
  \item 机器学习 — 用户-物品评分矩阵，大部分评分未知
\end{itemize}

\vspace{8pt}
\pause
\textbf{存储动机：} 对 $n=10^5$ 的 2D 网格问题，矩阵含 $10^{10}$ 个元素，但每行非零元仅 5 个。
若按稠密格式存储需要 $\sim$80 GB，稀疏格式仅需 $\sim$16 MB — \textbf{节省约三个数量级}。
\end{frame}

% ==================== frame 3: 可视化 ====================
\begin{frame}[shrink]{稀疏矩阵的可视化}
\framesubtitle{非零元分布模式}

\begin{center}
\begin{tikzpicture}[scale=0.85]
  % 5-pt stencil matrix pattern (tridiagonal blocks)
  \fill[blue!30] (0,0) rectangle (6,6);
  \fill[white] (0.3,0.3) rectangle (5.7,5.7);

  % main tridiagonal region
  \foreach \j in {0,1,2,3,4,5} {
    \fill[blue!60] (0.3+\j*0.9, 0.3+\j*0.9) rectangle (1.2+\j*0.9, 1.2+\j*0.9);
    \fill[blue!60] (0.3+\j*0.9, 1.2+\j*0.9) rectangle (1.2+\j*0.9, 2.1+\j*0.9);
    \fill[blue!60] (1.2+\j*0.9, 0.3+\j*0.9) rectangle (2.1+\j*0.9, 1.2+\j*0.9);
  }

  \draw[thick] (0,0) rectangle (6,6);
  \node[font=\footnotesize] at (3,-0.5) {5-点差分模板 — 块三对角结构};

  % arrows point to block detail
  \draw[->, thick] (7.0, 3.0) -- (8.5, 3.0);

  % detail: 5-point stencil connectivity
  \begin{scope}[shift={(10.0, 1.0)}, scale=1.3]
    \draw[gray, thin] (0,0) grid (3,3);
    \fill[blue!60] (1,1) circle (0.12); % center
    \fill[blue!40] (0,1) circle (0.1);  % left
    \fill[blue!40] (2,1) circle (0.1);  % right
    \fill[blue!40] (1,0) circle (0.1);  % bottom
    \fill[blue!40] (1,2) circle (0.1);  % top
    \draw[->, thick, red] (1,1) -- (0,1);
    \draw[->, thick, red] (1,1) -- (2,1);
    \draw[->, thick, red] (1,1) -- (1,0);
    \draw[->, thick, red] (1,1) -- (1,2);
    \node[font=\footnotesize] at (1.5,-0.4) {$2N \times 2N$ 网格};
  \end{scope}
\end{tikzpicture}
\end{center}

\vspace{4pt}
\textbf{稀疏模式决定存储格式选择：} 规则网格用 DIA/ELL，非规则图用 CSR/CSC。
\end{frame}

% ==================== frame 4: 常见存储格式概览 ====================
\begin{frame}[shrink]{常见稀疏矩阵存储格式}
\framesubtitle{六大经典格式及其适用场景}

\footnotesize
\begin{center}
\renewcommand{\arraystretch}{1.15}
\setlength{\tabcolsep}{3pt}
\begin{tabular}{@{}p{2.3cm} p{1.6cm} p{4.6cm} p{3.6cm}@{}}
\toprule
\textbf{格式} & \textbf{数组数} & \textbf{核心思想} & \textbf{适用场景} \\
\midrule
COO              & 3 & 直接存 (行, 列, 值) 三元组             & 构造期, I/O 格式 \\
\textbf{CSR (CRS)} & 3 & 行指针 + 列索引 + 值, 按行压缩       & \textbf{通用, SpMV 首选} \\
CSC              & 3 & 列指针 + 行索引 + 值, 按列压缩         & 列操作, 直接求解器 \\
DIA              & 2 & 存储各对角线                            & 规则网格 (每对角等长) \\
ELL (ELLPACK)    & 2 & 固定每行最大非零数, 补零                & 向量化, GPU 常用 \\
HYB              & 混合 & ELL + COO 组合                       & GPU, 平衡填充 \\
\bottomrule
\end{tabular}
\end{center}
\normalsize

\vspace{4pt}
\small
\textbf{核心权衡：} 存储紧凑性 vs 访问模式 vs 插入修改灵活性。

\textbf{CSR 优势：} SpMV 行访问连续, cache 友好; 压缩率高, 科学计算最通用格式。
\end{frame}

% ==================== frame 5: CRS 格式详解 ====================
\begin{frame}[shrink]{CRS (Compressed Row Storage) 格式详解}
\framesubtitle{三数组表示}

\textbf{给定稀疏矩阵 $A_{m \times n}$ 含 $N_{\!nz}$ 个非零元素，用三个数组表示：}

\begin{columns}[T]
\column{0.45\textwidth}
\[
A = \begin{pmatrix}
10 & 0  & 0  & 0  \\
0  & 20 & 0  & 30 \\
0  & 0  & 40 & 0  \\
50 & 0  & 0  & 60
\end{pmatrix}
\]

\column{0.55\textwidth}
\small
\textbf{val[ ]} — 非零元素值 (按行遍历)\\
\textbf{col\_ind[ ]} — 每个非零元的列索引\\
\textbf{row\_ptr[ ]} — 第 $i$ 行首元素在 val 中的偏移

\vspace{2pt}
\renewcommand{\arraystretch}{1.1}
\begin{tabular}{c|cccccc}
$i$ & 0 & 1 & 2 & 3 & 4 & 5 \\ \hline
val & 10 & 20 & 30 & 40 & 50 & 60 \\
col\_ind & 0 & 1 & 3 & 2 & 0 & 3 \\
\end{tabular}

\vspace{2pt}
\begin{tabular}{c|ccccc}
$i$ & 0 & 1 & 2 & 3 & 4 \\ \hline
row\_ptr & 0 & 1 & 3 & 4 & 6 \\
\end{tabular}
\end{columns}

\vspace{6pt}
\textbf{关键性质：}
\begin{itemize}
  \item row\_ptr[$i$] 到 row\_ptr[$i+1$]$-$1 为第 $i$ 行在 val/col\_ind 中的范围
  \item 第 $i$ 行非零元个数 $=$ row\_ptr[$i+1$] $-$ row\_ptr[$i$]；row\_ptr 长 $m+1$，末元素 $=N_{\!nz}$
\end{itemize}
\end{frame}

% ==================== frame 6: CRS 构造过程 ====================
\begin{frame}[shrink]{CRS 格式的构造过程}
\framesubtitle{从矩阵到三数组的变换}

\textbf{步骤：}
\begin{enumerate}
  \item \textbf{行扫描：} 按行遍历矩阵，收集每行非零元个数 $\to$ 计算 row\_ptr
  \item \textbf{偏移计算：} row\_ptr[$i+1$] = row\_ptr[$i$] + 第 $i$ 行非零元个数 (前缀和)
  \item \textbf{填充 val/col\_ind：} 再次按行遍历，将值和列索引填入对应位置
\end{enumerate}

\vspace{6pt}
\textbf{示例 ($4\times 4$ 矩阵, $N_{\!nz}=6$)：}

\small
\renewcommand{\arraystretch}{1.3}
\begin{tabular}{c|c|c|c}
\textbf{行} & \textbf{非零列} & \textbf{个数} & \textbf{row\_ptr 计算} \\
\hline
0 & col 0: 10                    & 1 & row\_ptr[0] = 0, row\_ptr[1] = 1 \\
1 & col 1: 20, col 3: 30         & 2 & row\_ptr[2] = 3 \\
2 & col 2: 40                    & 1 & row\_ptr[3] = 4 \\
3 & col 0: 50, col 3: 60         & 2 & row\_ptr[4] = 6 \\
\end{tabular}
\normalsize

\vspace{6pt}
\textbf{时间复杂度：} $O(m + N_{\!nz})$ — 两遍扫描，线性于矩阵规模。
空间复杂度：$O(N_{\!nz})$ — 比原始矩阵少存所有零元素。
\end{frame}

% ==================== frame 7: SpMV 算法 ====================
\begin{frame}[fragile]{CRS 核心操作：稀疏矩阵向量乘 (SpMV)}
\framesubtitle{$y \leftarrow A x$ — 迭代求解器最核心操作，占比超 80\%}

\small
\begin{lstlisting}[language=C,basicstyle=\ttfamily\footnotesize,belowskip=-4pt]
void spmv_csr(int m, const double *val, const int *col_ind,
              const int *row_ptr, const double *x, double *y)
{
  for (int i = 0; i < m; i++) {
    double sum = 0.0;
    for (int k = row_ptr[i]; k < row_ptr[i+1]; k++)
      sum += val[k] * x[col_ind[k]];
    y[i] = sum;
  }
}
\end{lstlisting}

\vspace{2pt}
\footnotesize
\textbf{访存模式：} val[k]、col\_ind[k] 连续读取 (stride-1)；$x$ 间接寻址 (随机访存，主瓶颈)；$y$ 连续写入。

\vspace{2pt}
\textbf{性能：} $\sim 2N_{\!nz}$ flops，算术强度 $\approx 0.17$ flops/byte — 典型的 memory-bound 操作。
\end{frame}

% ==================== frame 8: C 语言 CSR 数据结构 ====================
\begin{frame}[fragile]{C 语言中的 CSR 数据结构设计}

\begin{lstlisting}[language=C,basicstyle=\ttfamily\footnotesize]
typedef struct {
    int     m, n;       /* 行数、列数        */
    int     nnz;        /* 非零元总数        */
    double *val;        /* 非零元值 [nnz]    */
    int    *col_ind;    /* 列索引   [nnz]    */
    int    *row_ptr;    /* 行偏移   [m+1]    */
} csr_matrix;
\end{lstlisting}

\vspace{2pt}
\small
\textbf{设计：} \texttt{row\_ptr} 长度 $m+1$，末元素哨兵；堆分配，统一管理；内存 $8N_{\!nz} + 4(N_{\!nz}+m+1)$ 字节。
\end{frame}

% ==================== frame 9: COO → CSR 转换 ====================
\begin{frame}[fragile]{从 COO 格式构造 CSR}
\framesubtitle{两遍扫描算法 — 时间 $O(m+N_{\!nz})$}

\scriptsize
\begin{lstlisting}[language=C,basicstyle=\ttfamily\scriptsize,belowskip=-4pt]
csr->row_ptr = calloc(m + 1, sizeof(int));

/* Pass 1: count non-zeros per row */
for (int k = 0; k < nnz; k++)
    csr->row_ptr[coo_row[k] + 1]++;

/* Prefix sum to compute offsets */
for (int i = 1; i <= m; i++)
    csr->row_ptr[i] += csr->row_ptr[i - 1];

/* Pass 2: scatter values and column indices */
int *counter = calloc(m, sizeof(int));
for (int k = 0; k < nnz; k++) {
    int r = coo_row[k];
    int dest = csr->row_ptr[r] + counter[r]++;
    csr->val[dest]     = coo_val[k];
    csr->col_ind[dest] = coo_col[k];
}
free(counter);
\end{lstlisting}

\footnotesize
\textbf{要点：} 第 1 遍统计 $\to$ 前缀和得偏移 $\to$ 第 2 遍 counter 追踪写入位置。
\end{frame}

% ==================== frame 10: 前代/回代求解 ====================
\begin{frame}[fragile]{CSR 格式下的前代与回代}
\framesubtitle{三角稀疏线性系统的求解}

\footnotesize
\textbf{前代 $L\mathbf{x}=\mathbf{b}$ ($L$ 单位下三角)：}
\vspace{-4pt}
\begin{lstlisting}[language=C,basicstyle=\ttfamily\scriptsize,aboveskip=2pt,belowskip=2pt]
void csr_forward_sub(const csr_matrix *L,
                     const double *b, double *x) {
  for (int i = 0; i < L->m; i++) {
    double sum = b[i];
    for (int k = L->row_ptr[i]; k < L->row_ptr[i+1]; k++)
      if (L->col_ind[k] < i)
        sum -= L->val[k] * x[L->col_ind[k]];
    x[i] = sum;
  }
}
\end{lstlisting}

\vspace{2pt}
\textbf{回代 $U\mathbf{x}=\mathbf{y}$：} 从 $i=m-1$ 到 $0$，取 $j>i$ 的列。CSR 行遍历天然适合行优先的前代/回代。
\end{frame}

% ==================== frame 11: 矩阵转置 ====================
\begin{frame}[fragile]{CSR 格式下的转置操作}
\framesubtitle{CSR $\to$ CSC: 按行压缩的 $A^T$ 即按列压缩的 $A$}

\scriptsize
\begin{lstlisting}[language=C,basicstyle=\ttfamily\scriptsize,belowskip=-4pt]
AT->row_ptr = calloc(A->n + 1, sizeof(int));
for (int k = 0; k < A->nnz; k++)              /* count */
    AT->row_ptr[A->col_ind[k] + 1]++;
for (int i = 1; i <= A->n; i++)               /* prefix sum */
    AT->row_ptr[i] += AT->row_ptr[i - 1];

int *counter = calloc(A->n, sizeof(int));
for (int i = 0; i < A->m; i++)                /* scatter */
    for (int k = A->row_ptr[i]; k < A->row_ptr[i+1]; k++) {
        int j = A->col_ind[k];
        int dest = AT->row_ptr[j] + counter[j]++;
        AT->val[dest] = A->val[k];
        AT->col_ind[dest] = i;
    }
free(counter);
\end{lstlisting}

\vspace{2pt}
\footnotesize
同样使用计数 + 前缀和 + 散射三步框架，复杂度 $O(m+N_{\!nz})$。
\end{frame}

% ==================== frame 12: 性能优化 ====================
\begin{frame}[shrink]{CSR SpMV 性能优化}
\framesubtitle{内存墙 \textperiodcentered\ 向量化 \textperiodcentered\ 多线程}

\small
\textbf{核心矛盾：} SpMV 算术强度仅 $\sim$0.17 flops/byte，严重 memory-bound。

\vspace{4pt}
\begin{columns}[T]
\column{0.5\textwidth}
\textbf{存储格式层面}
\begin{itemize}
  \item SELL-C-$\sigma$: 行按长度分块，块内对齐 SIMD
  \item 分块 CSR (BCSR): 稠密小块 + 寄存器分块
  \item 混合格式: ELL 存规则部分 + COO 存异常行
  \item DIA: 规则网格问题效率最高
\end{itemize}

\column{0.5\textwidth}
\textbf{实现技术层面}
\begin{itemize}
  \item SIMD 向量化: sparse gather/scatter
  \item Cache 分块: $y$ 分块适配 cache，减少 $x$ 随机访存
  \item OpenMP: \texttt{\#pragma omp parallel for} 行间并行
  \item NUMA 感知: 矩阵和向量同节点
\end{itemize}
\end{columns}

\vspace{4pt}
\textbf{经验：} $n>10^5$ 时格式选择 $>$ 微调代码。CSR-SpMV $\sim$5--15 GFLOPS (memory-bound)，
而稠密 GEMM 可达 50+ GFLOPS。
\end{frame}

% ==================== frame 13: 完整示例 ====================
\begin{frame}[fragile]{完整示例 —— 从矩阵到 CSR 再到 SpMV}
\framesubtitle{端到端 C 代码}

\footnotesize
\begin{lstlisting}[language=C,basicstyle=\ttfamily\footnotesize,belowskip=-4pt]
int m = 4, n = 4, nnz = 6;
double val[] = {10, 20, 30, 40, 50, 60};
int    col[] = {0,  1,  3,  2,  0,  3 };
int    row[] = {0,  1,  3,  4,  6};  // m+1

csr_matrix A = {m, n, nnz, val, col, row};
double x[] = {1, 2, 3, 4}, y[4];
spmv_csr(A.m, A.val, A.col_ind, A.row_ptr, x, y);
// y = [10, 160, 120, 290]
\end{lstlisting}

\vspace{4pt}
\small
\textbf{验证：} $A=\left(\begin{smallmatrix}10&0&0&0\\0&20&0&30\\0&0&40&0\\50&0&0&60\end{smallmatrix}\right)$,
$x=(1,2,3,4)^T$,
$y=Ax=(10,\,20\cdot2+30\cdot4,\;40\cdot3,\;50\cdot1+60\cdot4)^T=(10,160,120,290)^T$.
\end{frame}

% ==================== frame 14: 性能对比 ====================
\begin{frame}[shrink]{存储格式性能对比 (SpMV)}
\framesubtitle{不同格式在同一矩阵上的表现}

\footnotesize
\begin{center}
\renewcommand{\arraystretch}{1.1}
\begin{tabular}{@{}l c c c c c@{}}
\toprule
\textbf{格式} & \textbf{存储(MB)} & \textbf{GFLOPS} & \textbf{带宽利用率} & \textbf{适用模式} & \textbf{GPU} \\
\midrule
Dense      & 80.0   & 52.0  & 95\%  & —        & 是 \\
\textbf{CSR} & \textbf{0.38} & \textbf{8.5} & \textbf{42\%} & \textbf{任意} & 一般 \\
CSC        & 0.38   & 7.8   & 38\%  & 任意     & 一般 \\
DIA        & 0.21   & 42.0  & 85\%  & 规则网格 & 是 \\
ELL        & 0.48   & 18.2  & 60\%  & 行等宽   & 是 \\
SELL-C-$\sigma$ & 0.40 & 22.5 & 72\%  & 较规则   & 是 \\
HYB        & 0.45   & 35.1  & 80\%  & 较规则   & 是 \\
\bottomrule
\end{tabular}
\end{center}

\vspace{2pt}
{\footnotesize \textit{条件：} $n=10^5$ 五点差分, $\approx$5e5 非零元, x86\_64 单核双精度.}

\vspace{2pt}
\small
\textbf{观察：} CSR 最通用但带宽利用率仅 40\% (随机 $x$ 访存)；DIA 在规则网格上可达 80\%；
GPU 上 ELL/HYB 显著优于 CSR (合并访存友好)。
\end{frame}

% ==================== frame 15: 总结 ====================
\begin{frame}[shrink]{总结}
\framesubtitle{关键要点与 C 实现注意事项}

\footnotesize
\begin{columns}[T]
\column{0.5\textwidth}
\textbf{核心概念}
\begin{enumerate}\itemsep1pt
  \item \textbf{稀疏性利用：} 只存非零元, 从 $O(n^2)$ 降至 $O(N_{\!nz})$
  \item \textbf{CRS：} val / col\_ind / row\_ptr 三数组, 按行存
  \item \textbf{构造：} 两遍扫描 $O(m+N_{\!nz})$, 前缀和为核心
  \item \textbf{SpMV：} 外层行循环 + 内层稀疏点积, memory-bound
  \item \textbf{格式选型：} CSR 通用, 规则网格 DIA, GPU 用 ELL/HYB
\end{enumerate}

\column{0.5\textwidth}
\textbf{C 实现注意事项}
\begin{itemize}\itemsep1pt
  \item \texttt{row\_ptr} 用 \texttt{int}, $n>2^{31}{-}1$ 改 \texttt{long long}
  \item 用 \texttt{calloc} 初始化 \texttt{row\_ptr}, 避免手动置零
  \item \texttt{counter} 数组追踪每行写入位置
  \item SpMV 内层循环 \texttt{restrict} 指针暗示无别名
  \item 反复调用的 SpMV: 将 $x$ 分块改善 cache 重用
  \item 构造期 COO 逐步添加, 完成后转 CSR
\end{itemize}
\end{columns}
\normalsize

\vspace{6pt}
\begin{center}
\textbf{稀疏矩阵计算的核心：用存储换访存, 用格式匹配结构。}
\end{center}
\end{frame}

\end{document}
